% !TEX root = ../../../main.tex
% 所属：第二部分：多变量分析学 / 第2章：多元函数微分学
% 节：2.5 无条件极值
% 正文按原文件顺序保存；由 ../chapter2.tex 调用。
\section{无条件极值}

\begin{definition}
设 $D\subset\mathbb{R}^n$ 为开区域，$f(\vect x)$ 为定义在 $D$ 上的函数，
$\vect x_0=(x_1^0,x_2^0,\ldots,x_n^0)\in D$。若存在 $\vect x_0$ 的邻域
$O(\vect x_0,r)$，使得
\[
f(\vect x_0)\geqslant f(\vect x)\quad
\bigl(\text{或 }f(\vect x_0)\leqslant f(\vect x)\bigr),
\qquad \vect x\in O(\vect x_0,r),
\]
则称 $\vect x_0$ 为 $f$ 的\textbf{极大值点}（或\textbf{极小值点}）；相应地，称 $f(\vect x_0)$ 为
相应的\textbf{极大值}（或\textbf{极小值}）；极大值点与极小值点统称为\textbf{极值点}，
极大值与极小值统称为\textbf{极值}。
\end{definition}
类似于单变量函数中的Fermat引理,我们有以下的必要条件:


\begin{proposition}
	若 $f(x_1,\ldots,x_n)$ 在
$(x_1^0,\ldots,x_n^0)$ 处可导，则 $(x_1^0,\ldots,x_n^0)$ 为
$f(x_1,\ldots,x_n)$ 的一个极值点的必要条件为
\[
\nabla f(\vect{x}^{\,0})=\vect{0}.
\]
\end{proposition}
我们称 $\nabla f(\vect{x}^{\,0})=\vect{0}$ 的点为驻点。

\textbf{注：}\begin{enumerate}
	\item 若 $f(x_1,\ldots,x_n)$ 在 $\vect{x}^{\,0}$ 处不可导，也可能有极值点。
	\item 驻点未必是极值点。例如 $z=xy$，$\nabla f=(y,x)$，$(0,0)$ 为驻点，
	但不是极值点。
\end{enumerate}

\begin{proof}
考虑
\[
\varphi(x_i)=f(x_1^0,\ldots,x_i,\ldots,x_n^0),
\qquad i=1,\ldots,n.
\]
$\varphi(x_i)$ 可导。$x_i^0$ 对 $\varphi(x_i)$ 为极值点，因此
\[
\frac{\mathrm d\varphi}{\mathrm d x_i}(x_i^0)=0
\quad\Longleftrightarrow\quad
\frac{\partial f}{\partial x_i}(\vec{x}^{\,0})=0.
\]
故
\[
\frac{\partial f}{\partial x_i}(\vect{x}^{\,0})=0,
\qquad i=1,\ldots,n.
\]
\end{proof}
\newpage
以 $z=f(x,y)$ 为例，设 $(x_0,y_0)$ 为 $f(x,y)$ 的一个驻点，
而且 $z=f(x,y)$ 在 $(x_0,y_0)$ 的某邻域
$D((x_0,y_0),\delta)$ 上二阶连续可导。在什么条件下，
$(x_0,y_0)$ 为 $z=f(x,y)$ 的一个极值点？类似于单变量函数中的做法，我们可以使用Lagrange中值定理来处理这个问题

取 $(x_0+\Delta x,y_0+\Delta y)\in D((x_0,y_0),\delta)$，则
\[
\begin{aligned}
&f(x_0+\Delta x,y_0+\Delta y)-f(x_0,y_0)\\
={}&\frac{\partial f}{\partial x}(x_0,y_0)\Delta x
+\frac{\partial f}{\partial y}(x_0,y_0)\Delta y\\
&+\frac12\left(
\frac{\partial^2 f}{\partial x^2}(x_0,y_0)\Delta x^2
+2\frac{\partial^2 f}{\partial x\partial y}(x_0,y_0)
\Delta x\Delta y
+\frac{\partial^2 f}{\partial y^2}(x_0,y_0)\Delta y^2
\right)+h\\
={}&\frac12\left(
\frac{\partial^2 f}{\partial x^2}(x_0,y_0)\Delta x^2
+2\frac{\partial^2 f}{\partial x\partial y}(x_0,y_0)
\Delta x\Delta y
+\frac{\partial^2 f}{\partial y^2}(x_0,y_0)\Delta y^2
\right)+h\\
={}&\frac12\left(A\Delta x^2+2B\Delta x\Delta y+C\Delta y^2\right)+h\\
={}&\frac12
\begin{pmatrix}\Delta x&\Delta y\end{pmatrix}
\begin{pmatrix}A&B\\B&C\end{pmatrix}
\begin{pmatrix}\Delta x\\\Delta y\end{pmatrix}+h,
\end{aligned}
\]
其中
\[
A=\frac{\partial^2 f}{\partial x^2}(x_0,y_0),
\qquad
B=\frac{\partial^2 f}{\partial x\partial y}(x_0,y_0),
\qquad
C=\frac{\partial^2 f}{\partial y^2}(x_0,y_0).
\]

\begin{enumerate}
\item 若 $\begin{pmatrix}A&B\\B&C\end{pmatrix}$ 正定，则当
$|\Delta x|,|\Delta y|$ 充分小时，
$f(x_0+\Delta x,y_0+\Delta y)>f(x_0,y_0)$，
$(x_0,y_0)$ 是严格极小值点。这等价于
\[
A>0\quad\text{且}\quad H=AC-B^2>0.
\]
\item 若 $\begin{pmatrix}A&B\\B&C\end{pmatrix}$ 负定，即
$H=AC-B^2>0$ 且 $A<0$，则 $(x_0,y_0)$ 是严格极大值点。
\item 若 $H<0$，则不是极值点。
\item 若 $H=0$，则不能确定。
\end{enumerate}

\begin{theorem}
设 $n$ 元函数 $f(x)$ 在
$x_0=(x_1^0,x_2^0,\ldots,x_n^0)$ 附近具有二阶连续偏导数，
且 $x_0$ 为 $f(x)$ 的驻点。那么当二次型
\[
g(\xi)=\sum_{i,j=1}^{n}f_{x_i x_j}(x_0)\xi_i\xi_j
\]
正定时，$f(x_0)$ 为极小值；当 $g(\xi)$ 负定时，
$f(x_0)$ 为极大值；当 $g(\xi)$ 不定时，$f(x_0)$ 不是极值。
\end{theorem}

记 $a_{ij}=f_{x_i x_j}(x_0)$，并记
\[
A_k=
\begin{pmatrix}
a_{11}&a_{12}&\cdots&a_{1k}\\
a_{21}&a_{22}&\cdots&a_{2k}\\
\vdots&\vdots&\ddots&\vdots\\
a_{k1}&a_{k2}&\cdots&a_{kk}
\end{pmatrix}.
\]
它称为 $f$ 的 $k$ 阶 Hessian 矩阵。由代数学知识即可得到：

\begin{corollary}
若 $\det A_k>0\ (k=1,2,\ldots,n)$，则二次型 $g(\xi)$ 是正定的，
此时 $f(x_0)$ 为极小值；若
$(-1)^k\det A_k>0\ (k=1,2,\ldots,n)$，
则二次型 $g(\xi)$ 是负定的，此时 $f(x_0)$ 为极大值。
\end{corollary}

若 $g(\xi)$ 不定，则 $\vect x_0$ 不是极值点；若 $g(\xi)$ 半正定，
则不能确定。

\subsection*{最小二乘法}

对一组数据
\[
(x_1,y_1),\ldots,(x_n,y_n)
\]
分布在 $y=ax+b$ 附近，使误差最小。令
\[
E=\sum_{k=1}^{n}|ax_k+b-y_k|^2,
\qquad
S(a,b)=\sum_{k=1}^{n}(ax_k+b-y_k)^2.
\]
若 $E(a,b)$ 取得最小值，则
\[
\begin{aligned}
\frac{\partial S}{\partial a}
&=\sum_{k=1}^{n}2x_k(ax_k+b-y_k)=0,\\
\frac{\partial S}{\partial b}
&=\sum_{k=1}^{n}2(ax_k+b-y_k)
 =2a\sum_{k=1}^{n}x_k+2nb-2\sum_{k=1}^{n}y_k=0.
\end{aligned}
\]
因而
\[
\begin{cases}
a\displaystyle\sum_{k=1}^{n}x_k^2
 +b\displaystyle\sum_{k=1}^{n}x_k
 =\displaystyle\sum_{k=1}^{n}x_ky_k,\\[6pt]
a\displaystyle\sum_{k=1}^{n}x_k+nb
 =\displaystyle\sum_{k=1}^{n}y_k.
\end{cases}
\]
解得
\[
\begin{aligned}
a&=
\frac{
n\displaystyle\sum_{k=1}^{n}x_ky_k
-\displaystyle\sum_{k=1}^{n}x_k\sum_{k=1}^{n}y_k
}{
n\displaystyle\sum_{k=1}^{n}x_k^2
-\left(\displaystyle\sum_{k=1}^{n}x_k\right)^2
},\\[8pt]
b&=
\frac{
\displaystyle\sum_{k=1}^{n}x_k^2\sum_{k=1}^{n}y_k
-\displaystyle\sum_{k=1}^{n}x_k\sum_{k=1}^{n}x_ky_k
}{
n\displaystyle\sum_{k=1}^{n}x_k^2
-\left(\displaystyle\sum_{k=1}^{n}x_k\right)^2
}.
\end{aligned}
\]

有时候，$(x_k,y_k)$ 不一定分布在直线附近，但
$(x_k,\ln y_k)$、$(x_k,e^{y_k})$、$(\ln x_k,y_k)$、
$(e^{x_k},y_k)$ 大致分布在直线附近，
就先把相应的变量变换再拟合。

